% Môn: Vật lý
% Lớp: 11
% Chương: Chương I. Dao động
% Bài: L11C1 Ôn Tập chương dao động · Bài toán xác định thời điểm và quãng đường.
% ===== Câu 1 =====
% ID: L11C1B1-01-DS
% Mức: TH
\dangbt{Bài toán xác định thời điểm và quãng đường.}
\begin{ex}
% ID: L11C1B1-01-DS
% Mức: TH
Vật dao động điều hòa biên độ $4\,\text{cm}$, chu kì $1\,\text{s}$. Xác định tính đúng sai.
\choiceTF
{\True Quãng đường trong một chu kì là $16\,\text{cm}$.}
{\True Quãng đường trong nửa chu kì là $8\,\text{cm}$.}
{\True Sau một chu kì, độ dịch chuyển bằng 0.}
{\True Tốc độ trung bình trong một chu kì bằng $16\,\text{cm}$ /s.}
\end{ex}

% ===== Câu 2 =====
% ID: L11C1B1-02-DS
% Mức: TH
\begin{ex}
% ID: L11C1B1-02-DS
% Mức: TH
Vật dao động điều hòa biên độ $4\,\text{cm}$, chu kì $1\,\text{s}$. Tính quãng đường và tốc độ trung bình trong một chu kì.
\choiceTF

\end{ex}

% ===== Câu 3 =====
% ID: L11C1B1-03-TLN
% Mức: VD
\begin{ex}
% ID: L11C1B1-03-TLN
% Mức: VD
Cho một chất điểm dao động điều hòa quanh vị trí cân bằng O. Li độ biến thiên theo thời gian như mô tả hình bên dưới. Quỹ đạo dao động có độ dài bằng bao nhiêu mét? (Kết quả đến 2 chữ số sau dấu phẩy thập phân.)
\shortans{0,08}
\loigiai{
: Biên độ dao động là: $A$ = 4(cm)= 0,04(m) Chiều dài quỹ đạo dao động là: $L=2A=0,08(m)$
}
\end{ex}

% ===== Câu 4 =====
% ID: L11C1B1-04-TN
% Mức: NB
\begin{ex}
% ID: L11C1B1-04-TN
% Mức: NB
Một chất điểm dao động điều hòa có quỹ đạo là một đoạn thẳng dài $30$ cm. Biên độ dao động của chất điểm là
\choice
{$30$ cm}
{\True $15$ cm}
{$-15$ cm}
{$7,5\mathrm{~cm}$}
\loigiai{
Biên độ dao động là độ lệch cực đại của vật so với vị trí cân bằng, bằng một nửa quỹ đạo dao động: $A = \frac{30}{2} = 15 \mathrm{~cm}$
}
\end{ex}

% ===== Câu 5 =====
% ID: L11C1B1-05-DS
% Mức: VD
\begin{ex}
% ID: L11C1B1-05-DS
% Mức: VD
Vật dao động điều hòa biên độ $6\,\text{cm}$, chu kì $2\,\text{s}$. Xác định tính đúng sai.
\choiceTF
{\True Quãng đường trong một chu kì là $24\,\text{cm}$.}
{\True Quãng đường trong nửa chu kì là $12\,\text{cm}$.}
{\True Sau một chu kì, độ dịch chuyển bằng 0.}
{\True Tốc độ trung bình trong một chu kì bằng $12\,\text{cm}$ /s.}
\end{ex}

% ===== Câu 6 =====
% ID: L11C1B1-06-DS
% Mức: VD
\begin{ex}
% ID: L11C1B1-06-DS
% Mức: VD
Vật dao động điều hòa biên độ $6\,\text{cm}$, chu kì $2\,\text{s}$. Tính quãng đường và tốc độ trung bình trong một chu kì.
\choiceTF

\end{ex}

% ===== Câu 7 =====
% ID: L11C1B1-07-TLN
% Mức: VD
\begin{ex}
% ID: L11C1B1-07-TLN
% Mức: VD
Một vật nhỏ dao động điều hòa với biên độ $A = 5\text{ cm}$. Quãng đường vật đi được sau 2 dao động toàn phần (2 chu kì) theo cm :
\shortans{40}
\loigiai{
Trong một dao động toàn phần (1 chu kì $T$), quãng đường vật luôn đi được là $4A$. 
Sau 2 dao động toàn phần, quãng đường vật đi được là: $S = 2 \cdot 4A = 8A = 8 \cdot 5 = 40\text{ cm}$.
}
\end{ex}

% ===== Câu 8 =====
% ID: L11C1B1-08-TN
% Mức: VD
\begin{ex}
% ID: L11C1B1-08-TN
% Mức: VD
Một chất điểm dao động điều hòa với chu kì $T$. Trong khoảng thời gian ngắn nhất khi đi từ vị trí biên có li độ $x_1 = A$ đến vị trí $x_2 = A/2$, chất điểm có tốc độ trung bình là:
\choice
{$\dfrac{A}{T}$}
{$\dfrac{2A}{T}$}
{\True $\dfrac{3A}{T}$}
{$\dfrac{6A}{T}$}
\loigiai{
Chu kì dao động là $T$. Vị trí ban đầu là $x_1 = A$ (biên dương). Vị trí đích là $x_2 = A/2$. Quãng đường vật đi được trong khoảng thời gian ngắn nhất từ $x=A$ đến $x=A/2$ là $\Delta S = \left|A - \dfrac{A}{2}\right| = \dfrac{A}{2}$. Thời gian ngắn nhất để vật đi từ $x=A$ đến $x=A/2$ là $\Delta t = \dfrac{T}{6}$ (vì từ biên $A$ đến $A/2$ tương ứng góc pha $60^\circ = \dfrac{T}{6}$). Tốc độ trung bình là $v_{tb} = \dfrac{\Delta S}{\Delta t} = \dfrac{A/2}{T/6} = \dfrac{3A}{T}$.
}
\end{ex}

% ===== Câu 9 =====
% ID: L11C1B1-09-DS
% Mức: VD
\begin{ex}
% ID: L11C1B1-09-DS
% Mức: VD
Vật dao động điều hòa biên độ $8\,\text{cm}$, chu kì $4\,\text{s}$. Xác định tính đúng sai.
\choiceTF
{\True Quãng đường trong một chu kì là $32\,\text{cm}$.}
{\True Quãng đường trong nửa chu kì là $16\,\text{cm}$.}
{\True Sau một chu kì, độ dịch chuyển bằng 0.}
{\True Tốc độ trung bình trong một chu kì bằng $8\,\text{cm}$ /s.}
\end{ex}

% ===== Câu 10 =====
% ID: L11C1B1-10-DS
% Mức: VD
\begin{ex}
% ID: L11C1B1-10-DS
% Mức: VD
Vật dao động điều hòa biên độ $8\,\text{cm}$, chu kì $4\,\text{s}$. Tính quãng đường và tốc độ trung bình trong một chu kì.
\choiceTF

\end{ex}

% ===== Câu 11 =====
% ID: L11C1B1-11-TLN
% Mức: VDC
\begin{ex}
% ID: L11C1B1-11-TLN
% Mức: VDC
Một cái chốt có kích thước đáng kể được gắn lên mép một cái đĩa bán kính $r=8cm$ như hình vẽ. Đĩa quay với tốc độ góc không đổi $\omega =\dfrac{\pi }{4}rad/s.$ Một chùm tia sáng song song nằm ngang tạo ra bóng của cái chốt lên màn. Tại thời điểm ban đầu, cái chốt ở $P$ và bóng của nó lên màn tại $S$. Xác định bóng của cái chốt trên màn tại thời điểm $t=0,3s$ (kết quả làm tròn đến chữ số thập phân thứ 2 sau dấu phẩy)
\shortans{1,87}
\loigiai{
Phương trình dao động của chất điểm $x = 8\sin\left(\dfrac{\pi}{4}t + \varphi_0\right) \mathrm{cm}$. Tại $t = 0$ thì $x = 0$ → $\varphi_0 = 0 \, \mathrm{rad}$ → $x = 8\sin\left(\dfrac{\pi}{4}t\right) \mathrm{cm}$ (*). Thay $t = 0,3 \, \mathrm{s}$ vào (*) ta có $x \approx 1,86756 \, \mathrm{cm}$.
}
\end{ex}

% ===== Câu 12 =====
% ID: L11C1B1-12-TN
% Mức: VD
\begin{ex}
% ID: L11C1B1-12-TN
% Mức: VD
Một chất điểm dao động điều hòa trong 10 dao động toàn phần chất điểm đi được quãng đường dài $120\,\text{cm}$. Quỹ đạo dao động của vật có chiều dài là:
\choice
{\True $6\,\text{cm}$}
{$12\,\text{cm}$}
{$3\,\text{cm}$}
{$9\,\text{cm}$}
\loigiai{
Quỹ đạo dao động của vật có chiều dài là $2A=6$ cm.
  Còn quãng đường của 1 dao động là $4$ A.
}
\end{ex}

% ===== Câu 15 =====
% ID: L11C1B1-13-TLN
% Mức: VDC
\begin{ex}
% ID: L11C1B1-13-TLN
% Mức: VDC
Cho hai vật dao động điều hòa có phương trình lần lượt là $x_1 = 5\cos \left( 2\pi t + \dfrac{\pi}{3} \right)\,(\text{cm})$ và $x_2 = 7\cos \left( 2\pi t - \dfrac{\pi}{3} \right)\,(\text{cm})$. Khi li độ của dao động 1 đạt đến giá trị cực đại thì sau thời gian ngắn nhất là bao lâu li độ của dao động 2 đạt đến cực đại?
\shortans{0,25}
\loigiai{
Khi li độ của dao động 1 đạt đến giá trị cực đại thì sau thời gian ngắn nhất li độ của dao động 2 đạt đến cực đại là $\Delta t=\dfrac{\Delta \varphi }{\omega }=\dfrac{1}{4}$ s
}
\end{ex}

% ===== Câu 16 =====
% ID: L11C1B1-14-TN
% Mức: VD
\begin{ex}
% ID: L11C1B1-14-TN
% Mức: VD
Một chất điểm dao động điều hòa trong 10 dao động toàn phần chất điểm đi được quãng đường dài $120\,\text{cm}$. Quỹ đạo dao động của vật có chiều dài là:
\choice
{\True $6\,\text{cm}$}
{$12\,\text{cm}$}
{$3\,\text{cm}$}
{$9\,\text{cm}$}
\loigiai{
Quỹ đạo dao động của vật có chiều dài là $2A=6$ cm.
  Còn quãng đường của 1 dao động là $4$ A.
}
\end{ex}

% ===== Câu 19 =====
% ID: L11C1B1-15-TLN
% Mức: VDC
\begin{ex}
% ID: L11C1B1-15-TLN
% Mức: VDC
Vật dao động điều hòa theo phương trình $x = -5\cos(10\pi t)$ (cm). Thời gian vật đi quãng đường dài $12,5$ cm kể từ lúc bắt đầu chuyển động là:
\shortans{B}
\loigiai{
• Từ $x = -5$ đến $x = 0$: quãng đường $5$ cm.
  
 • Từ $x = 0$ đến $x = +5$: quãng đường $5$ cm.
  
 • Từ $x = +5$ đến $x = +2,5$: quãng đường $2,5$ cm.
  
  Thời gian đi từ $x = -5$ đến $x = +5$ là $\dfrac{T}{2} = \dfrac{1}{10}$ s.
Vật quay lại $x = +2,5$ nên ta cần tính thời gian từ biên về vị trí cách $2,5\,\text{cm}$:
  $x = A\cos(\omega t) = 2,5 \Rightarrow \cos(\omega t) = \dfrac{1}{2} \Rightarrow \omega t = \dfrac{\pi}{3} \Rightarrow t = \dfrac{1}{30}$
  Tổng thời gian: $\dfrac{1}{10} + \dfrac{1}{30} = \boxed{\dfrac{2}{15}}\,\text{s}$
}
\end{ex}

% ===== Câu 20 =====
% ID: L11C1B1-16-TN
% Mức: VD
\begin{ex}
% ID: L11C1B1-16-TN
% Mức: VD
Một chất điểm dao động điều hòa trong 10 dao động toàn phần chất điểm đi được quãng đường dài $120\,\text{cm}$. Quỹ đạo dao động của vật có chiều dài là:
\choice
{\True $6\,\text{cm}$}
{$12\,\text{cm}$}
{$3\,\text{cm}$}
{$9\,\text{cm}$}
\loigiai{
Quỹ đạo dao động của vật có chiều dài là $2A=6$ cm.
  Còn quãng đường của 1 dao động là $4$ A.
}
\end{ex}

% ===== Câu 23 =====
% ID: L11C1B1-17-TLN
% Mức: VDC
\begin{ex}
% ID: L11C1B1-17-TLN
% Mức: VDC
Vật dao động điều hòa theo phương trình $x = -5\cos(10\pi t)$ (cm). Thời gian vật đi quãng đường dài $12,5$ cm kể từ lúc bắt đầu chuyển động là:
\shortans{B}
\loigiai{
• Từ $x = -5$ đến $x = 0$: quãng đường $5$ cm.
  
 • Từ $x = 0$ đến $x = +5$: quãng đường $5$ cm.
  
 • Từ $x = +5$ đến $x = +2,5$: quãng đường $2,5$ cm.
  
  Thời gian đi từ $x = -5$ đến $x = +5$ là $\dfrac{T}{2} = \dfrac{1}{10}$ s.
Vật quay lại $x = +2,5$ nên ta cần tính thời gian từ biên về vị trí cách $2,5\,\text{cm}$:
  $x = A\cos(\omega t) = 2,5 \Rightarrow \cos(\omega t) = \dfrac{1}{2} \Rightarrow \omega t = \dfrac{\pi}{3} \Rightarrow t = \dfrac{1}{30}$
  Tổng thời gian: $\dfrac{1}{10} + \dfrac{1}{30} = \boxed{\dfrac{2}{15}}\,\text{s}$
}
\end{ex}

% ===== Câu 24 =====
% ID: L11C1B1-18-TN
% Mức: VD
\begin{ex}
% ID: L11C1B1-18-TN
% Mức: VD
Một chất điểm dao động điều hòa trong 10 dao động toàn phần chất điểm đi được quãng đường dài $120\,\text{cm}$. Quỹ đạo dao động của vật có chiều dài là:
\choice
{\True $6\,\text{cm}$}
{$12\,\text{cm}$}
{$3\,\text{cm}$}
{$9\,\text{cm}$}
\loigiai{
Quỹ đạo dao động của vật có chiều dài là $2A=6$ cm.
  Còn quãng đường của 1 dao động là $4$ A.
}
\end{ex}

% ===== Câu 27 =====
% ID: L11C1B1-19-TLN
% Mức: VDC
\begin{ex}
% ID: L11C1B1-19-TLN
% Mức: VDC
Vật dao động điều hòa theo phương trình $x = -5\cos(10\pi t)$ (cm). Thời gian vật đi quãng đường dài $12,5$ cm kể từ lúc bắt đầu chuyển động là:
\shortans{B}
\loigiai{
• Từ $x = -5$ đến $x = 0$: quãng đường $5$ cm.
  
 • Từ $x = 0$ đến $x = +5$: quãng đường $5$ cm.
  
 • Từ $x = +5$ đến $x = +2,5$: quãng đường $2,5$ cm.
  
  Thời gian đi từ $x = -5$ đến $x = +5$ là $\dfrac{T}{2} = \dfrac{1}{10}$ s.
Vật quay lại $x = +2,5$ nên ta cần tính thời gian từ biên về vị trí cách $2,5\,\text{cm}$:
  $x = A\cos(\omega t) = 2,5 \Rightarrow \cos(\omega t) = \dfrac{1}{2} \Rightarrow \omega t = \dfrac{\pi}{3} \Rightarrow t = \dfrac{1}{30}$
  Tổng thời gian: $\dfrac{1}{10} + \dfrac{1}{30} = \boxed{\dfrac{2}{15}}\,\text{s}$
}
\end{ex}

% ===== Câu 28 =====
% ID: L11C1B1-20-TN
% Mức: VDC
\begin{ex}
% ID: L11C1B1-20-TN
% Mức: VDC
Một vật thực hiện dao động điều hòa theo phương trình $x = 6\cos(10\pi t)$ cm. Tốc độ trung bình kể từ khi vật ở vị trí cân bằng đang chuyển động theo chiều dương đến thời điểm đầu tiên vật có li độ $3\text{ cm}$ là:
\choice
{\True $1,8\text{ m/s}$}
{$0,9\text{ m/s}$}
{$2,4\text{ m/s}$}
{$1,2\text{ m/s}$}
\loigiai{
Phương trình dao động $x = 6\cos(10\pi t)$ cm. Biên độ $A = 6\text{ cm}$, tần số góc $\omega = 10\pi\text{ rad/s}$, chu kì $T = \dfrac{2\pi}{\omega} = \dfrac{2\pi}{10\pi} = 0,2\text{ s}$. Trạng thái ban đầu: vật ở vị trí cân bằng ($x = 0$) và đang chuyển động theo chiều dương, tương ứng pha $\varphi_0 = -\dfrac{\pi}{2}$. Vị trí đích: $x = 3\text{ cm} = \dfrac{A}{2}$. Khi xuất phát từ $x = 0$ theo chiều dương, vật đến $x = \dfrac{A}{2}$ lần đầu tiên sau thời gian $\Delta t = \dfrac{T}{12} = \dfrac{0,2}{12} = \dfrac{1}{60}\text{ s}$. Quãng đường đi được: $\Delta S = \dfrac{A}{2} = 3\text{ cm}$. Tốc độ trung bình: $v_{tb} = \dfrac{\Delta S}{\Delta t} = \dfrac{3\text{ cm}}{\dfrac{1}{60}\text{ s}} = 180\text{ cm/s} = 1,8\text{ m/s}$.
}
\end{ex}

% ===== Câu 31 =====
% ID: L11C1B1-21-TLN
% Mức: VDC
\begin{ex}
% ID: L11C1B1-21-TLN
% Mức: VDC
Vật dao động điều hòa theo phương trình $x = -5\cos(10\pi t)$ (cm). Thời gian vật đi quãng đường dài $12,5$ cm kể từ lúc bắt đầu chuyển động là:
\shortans{B}
\loigiai{
• Từ $x = -5$ đến $x = 0$: quãng đường $5$ cm.
  
 • Từ $x = 0$ đến $x = +5$: quãng đường $5$ cm.
  
 • Từ $x = +5$ đến $x = +2,5$: quãng đường $2,5$ cm.
  
  Thời gian đi từ $x = -5$ đến $x = +5$ là $\dfrac{T}{2} = \dfrac{1}{10}$ s.
Vật quay lại $x = +2,5$ nên ta cần tính thời gian từ biên về vị trí cách $2,5\,\text{cm}$:
  $x = A\cos(\omega t) = 2,5 \Rightarrow \cos(\omega t) = \dfrac{1}{2} \Rightarrow \omega t = \dfrac{\pi}{3} \Rightarrow t = \dfrac{1}{30}$
  Tổng thời gian: $\dfrac{1}{10} + \dfrac{1}{30} = \boxed{\dfrac{2}{15}}\,\text{s}$
}
\end{ex}

% ===== Câu 32 =====
% ID: L11C1B1-22-TN
% Mức: VDC
\begin{ex}
% ID: L11C1B1-22-TN
% Mức: VDC
Một con lắc lò xo có độ cứng $1\text{N/m}$, vật nặng có khối lượng $100\text{g}$ dao động điều hòa theo phương ngang, trong quá trình dao động, vận tốc có độ lớn cực đại $6\pi \text{ cm/s}$, lấy $\pi^2 = 10$. Thời gian ngắn nhất vật đi từ vị trí $x = 6\text{ cm}$ đến vị trí $-3\sqrt{3}\text{ (cm)}$ là:
\choice
{$0,333$}
{$0,667$}
{$0,167$}
{\True $0,833$}
\loigiai{
Khối lượng $m = 100\text{g} = 0,1\text{ kg}$. Độ cứng $k = 1\text{ N/m}$. Vận tốc cực đại $v_{max} = 6\pi \text{ cm/s}$. Tần số góc của dao động: $\omega = \sqrt{\dfrac{k}{m}} = \sqrt{\dfrac{1}{0,1}} = \sqrt{10}\text{ rad/s}$. Vì $\pi^2 = 10$, ta có $\omega = \sqrt{\pi^2} = \pi\text{ rad/s}$. Biên độ dao động: $A = \dfrac{v_{max}}{\omega} = \dfrac{6\pi}{\pi} = 6\text{ cm}$. Chu kì dao động: $T = \dfrac{2\pi}{\omega} = \dfrac{2\pi}{\pi} = 2\text{ s}$. Vị trí ban đầu $x_1 = 6\text{ cm}$. Đây chính là biên dương ($x_1 = A$). Vị trí đích $x_2 = -3\sqrt{3}\text{ cm}$. Ta có $-3\sqrt{3} = -6\dfrac{\sqrt{3}}{2} = -A\dfrac{\sqrt{3}}{2}$. Thời gian ngắn nhất để vật đi từ $x=A$ đến $x=-A\dfrac{\sqrt{3}}{2}$: 
 • Từ $x=A$ đến vị trí cân bằng $x=0$: thời gian $t_1 = T/4$. 
 • Từ $x=0$ đến $x=-A\dfrac{\sqrt{3}}{2}$: thời gian $t_2 = T/6$. Tổng thời gian $\Delta t = t_1 + t_2 = T/4 + T/6 = \dfrac{3T}{12} + \dfrac{2T}{12} = \dfrac{5T}{12}$. $\Delta t = \dfrac{5 \times 2}{12} = \dfrac{10}{12} = \dfrac{5}{6}\text{ s}$. Giá trị xấp xỉ: $\dfrac{5}{6} \approx 0,8333\text{ s}$.
}
\end{ex}

% ===== Câu 35 =====
% ID: L11C1B1-23-TN
% Mức: VDC
\begin{ex}
% ID: L11C1B1-23-TN
% Mức: VDC
Một con lắc lò xo nằm ngang dao động điều hòa với biên độ $10\text{cm}$, chu kì $1\text{s}$. Trong một chu kì, quãng thời gian mà khoảng cách từ vật tới vị trí cân bằng lớn hơn $5\sqrt{2}\text{ cm}$ là:
\choice
{$\dfrac{1}{3}\text{ s}$}
{\True $0,5\text{ s}$}
{$\dfrac{1}{6}\text{ s}$}
{$\dfrac{1}{9}\text{ s}$}
\loigiai{
Biên độ $A = 10\text{ cm}$. Chu kì $T = 1\text{ s}$. Khoảng cách từ vật tới vị trí cân bằng lớn hơn $5\sqrt{2}\text{ cm}$ có nghĩa là $|x| {\gt } 5\sqrt{2}\text{ cm}$. Ta có $5\sqrt{2}\text{ cm} = \dfrac{10\sqrt{2}}{2} = A\dfrac{\sqrt{2}}{2}$. Thời gian vật có li độ thỏa mãn $|x| {\gt } A\dfrac{\sqrt{2}}{2}$ trong một chu kì là: Thời gian vật đi từ $x=A\dfrac{\sqrt{2}}{2}$ đến $x=A$ là $T/8$. Trong một chu kì, có $4$ lần vật đi qua vùng này (từ $A\dfrac{\sqrt{2}}{2} \to A$, từ $-A\dfrac{\sqrt{2}}{2} \to -A$ và các chiều ngược lại). Tổng thời gian là $\Delta t = 4 \times (T/8) = T/2$. $\Delta t = 1/2 = 0,5\text{ s}$.
}
\end{ex}

% ===== Câu 39 =====
% ID: L11C1B1-24-TN
% Mức: VDC
\begin{ex}
% ID: L11C1B1-24-TN
% Mức: VDC
Vật dao động điều hòa theo phương trình $x = \cos(\pi t - \pi/3)$ (dm). Thời gian vật đi được quãng đường $S = 5\text{ cm}$ kể từ thời điểm ban đầu ($t = 0$) là:
\choice
{\True $\dfrac{1}{3}$ s}
{$\dfrac{1}{6}$ s}
{$\dfrac{1}{4}$ s}
{$\dfrac{1}{2}$ s}
\loigiai{
Phương trình dao động $x = \cos(\pi t - \pi/3)$ dm. Biên độ $A = 1\text{ dm} = 10\text{ cm}$. Tần số góc $\omega = \pi\text{ rad/s}$. Chu kì $T = \dfrac{2\pi}{\omega} = \dfrac{2\pi}{\pi} = 2\text{ s}$. Tại $t = 0$: $x(0) = \cos(-\pi/3) = \dfrac{1}{2}\text{ dm} = 5\text{ cm} = \dfrac{A}{2}$. Vận tốc $v(0) = -\omega A\sin(-\pi/3) = -\pi \cdot 1 \cdot \left(-\dfrac{\sqrt{3}}{2}\right) = \dfrac{\pi\sqrt{3}}{2}\text{ dm/s} \gt  0$, vật đang chuyển động theo chiều dương. Vật xuất phát từ $x = A/2$ đi theo chiều dương đến biên $x = A$, quãng đường đúng bằng $A - A/2 = 5\text{ cm}$. Trên đường tròn pha, góc quét từ $\varphi = -\pi/3$ đến $\varphi = 0$ là $\pi/3$, tương ứng thời gian $\Delta t = \dfrac{\pi/3}{\omega} = \dfrac{\pi/3}{\pi} = \dfrac{1}{3}\text{ s}$. Vậy đáp án là $\dfrac{1}{3}$ s.
}
\end{ex}

% ===== Câu 43 =====
% ID: L11C1B1-25-TN
% Mức: VDC
\begin{ex}
% ID: L11C1B1-25-TN
% Mức: VDC
Vật dao động điều hòa theo phương trình $x = \cos(\pi t - \pi/3)$ (dm). Thời gian vật đi được quãng đường $S = 5\text{ cm}$ kể từ thời điểm ban đầu ($t = 0$) là:
\choice
{\True $\dfrac{1}{3}$ s}
{$\dfrac{1}{6}$ s}
{$\dfrac{1}{4}$ s}
{$\dfrac{1}{2}$ s}
\loigiai{
Phương trình dao động $x = \cos(\pi t - \pi/3)$ dm. Biên độ $A = 1\text{ dm} = 10\text{ cm}$. Tần số góc $\omega = \pi\text{ rad/s}$. Chu kì $T = \dfrac{2\pi}{\omega} = \dfrac{2\pi}{\pi} = 2\text{ s}$. Tại $t = 0$: $x(0) = \cos(-\pi/3) = \dfrac{1}{2}\text{ dm} = 5\text{ cm} = \dfrac{A}{2}$. Vận tốc $v(0) = -\omega A\sin(-\pi/3) = -\pi \cdot 1 \cdot \left(-\dfrac{\sqrt{3}}{2}\right) = \dfrac{\pi\sqrt{3}}{2}\text{ dm/s} \gt  0$, vật đang chuyển động theo chiều dương. Vật xuất phát từ $x = A/2$ đi theo chiều dương đến biên $x = A$, quãng đường đúng bằng $A - A/2 = 5\text{ cm}$. Trên đường tròn pha, góc quét từ $\varphi = -\pi/3$ đến $\varphi = 0$ là $\pi/3$, tương ứng thời gian $\Delta t = \dfrac{\pi/3}{\omega} = \dfrac{\pi/3}{\pi} = \dfrac{1}{3}\text{ s}$. Vậy đáp án là $\dfrac{1}{3}$ s.
}
\end{ex}

% ===== Câu 47 =====
% ID: L11C1B1-26-TN
% Mức: VDC
\begin{ex}
% ID: L11C1B1-26-TN
% Mức: VDC
Vật dao động điều hòa theo phương trình $x = \cos(\pi t - \pi/3)$ (dm). Thời gian vật đi được quãng đường $S = 5\text{ cm}$ kể từ thời điểm ban đầu ($t = 0$) là:
\choice
{\True $\dfrac{1}{3}$ s}
{$\dfrac{1}{6}$ s}
{$\dfrac{1}{4}$ s}
{$\dfrac{1}{2}$ s}
\loigiai{
Phương trình dao động $x = \cos(\pi t - \pi/3)$ dm. Biên độ $A = 1\text{ dm} = 10\text{ cm}$. Tần số góc $\omega = \pi\text{ rad/s}$. Chu kì $T = \dfrac{2\pi}{\omega} = \dfrac{2\pi}{\pi} = 2\text{ s}$. Tại $t = 0$: $x(0) = \cos(-\pi/3) = \dfrac{1}{2}\text{ dm} = 5\text{ cm} = \dfrac{A}{2}$. Vận tốc $v(0) = -\omega A\sin(-\pi/3) = -\pi \cdot 1 \cdot \left(-\dfrac{\sqrt{3}}{2}\right) = \dfrac{\pi\sqrt{3}}{2}\text{ dm/s} \gt  0$, vật đang chuyển động theo chiều dương. Vật xuất phát từ $x = A/2$ đi theo chiều dương đến biên $x = A$, quãng đường đúng bằng $A - A/2 = 5\text{ cm}$. Trên đường tròn pha, góc quét từ $\varphi = -\pi/3$ đến $\varphi = 0$ là $\pi/3$, tương ứng thời gian $\Delta t = \dfrac{\pi/3}{\omega} = \dfrac{\pi/3}{\pi} = \dfrac{1}{3}\text{ s}$. Vậy đáp án là $\dfrac{1}{3}$ s.
}
\end{ex}

% ===== Câu 49 =====
% ID: L11C1B1-27-TN
% Mức: VDC
\begin{ex}
% ID: L11C1B1-27-TN
% Mức: VDC
Vật dao động điều hòa theo phương trình $x = \cos(\pi t - \pi/3)$ (dm). Thời gian vật đi được quãng đường $S = 5\text{ cm}$ kể từ thời điểm ban đầu ($t = 0$) là:
\choice
{\True $\dfrac{1}{3}$ s}
{$\dfrac{1}{6}$ s}
{$\dfrac{1}{4}$ s}
{$\dfrac{1}{2}$ s}
\loigiai{
Phương trình dao động $x = \cos(\pi t - \pi/3)$ dm. Biên độ $A = 1\text{ dm} = 10\text{ cm}$. Tần số góc $\omega = \pi\text{ rad/s}$. Chu kì $T = \dfrac{2\pi}{\omega} = \dfrac{2\pi}{\pi} = 2\text{ s}$. Tại $t = 0$: $x(0) = \cos(-\pi/3) = \dfrac{1}{2}\text{ dm} = 5\text{ cm} = \dfrac{A}{2}$. Vận tốc $v(0) = -\omega A\sin(-\pi/3) = -\pi \cdot 1 \cdot \left(-\dfrac{\sqrt{3}}{2}\right) = \dfrac{\pi\sqrt{3}}{2}\text{ dm/s} \gt  0$, vật đang chuyển động theo chiều dương. Vật xuất phát từ $x = A/2$ đi theo chiều dương đến biên $x = A$, quãng đường đúng bằng $A - A/2 = 5\text{ cm}$. Trên đường tròn pha, góc quét từ $\varphi = -\pi/3$ đến $\varphi = 0$ là $\pi/3$, tương ứng thời gian $\Delta t = \dfrac{\pi/3}{\omega} = \dfrac{\pi/3}{\pi} = \dfrac{1}{3}\text{ s}$. Vậy đáp án là $\dfrac{1}{3}$ s.
}
\end{ex}

% ===== Câu 51 =====
% ID: L11C1B1-28-TN
% Mức: VDC
\begin{ex}
% ID: L11C1B1-28-TN
% Mức: VDC
Một vật dao động điều hòa theo phương trình $x = 4\cos(8\pi t - 2\pi/3)$ cm. Thời gian vật đi được quãng đường $s = (2+2\sqrt{2})$ cm kể từ lúc vật bắt đầu dao động là:
\choice
{\True $\dfrac{5}{96}$ (s)}
{$\dfrac{29}{96}$ (s)}
{$\dfrac{1}{96}$ (s)}
{$\dfrac{1}{48}$ (s)}
\loigiai{
Pha ban đầu: $x(0) = 4\cos(-\tfrac{2\pi}{3}) = 4(-\tfrac{1}{2}) = -2$ cm. Vận tốc ban đầu: $v(0) = -A\omega\sin(-\tfrac{2\pi}{3}) = -4(8\pi)(-\tfrac{\sqrt{3}}{2}) = 16\pi\sqrt{3} \gt  0$. Vật đi theo chiều dương. Quãng đường cần đi: $s = (2+2\sqrt{2})$ cm. Bước 1: Vật đi từ $x_0 = -2$ cm đến $x=0$. Quãng đường $s_1 = |0 - (-2)| = 2$ cm. Góc quay tương ứng: từ $\theta_0 = -\tfrac{2\pi}{3}$ đến $\theta_1 = -\tfrac{\pi}{2}$. Thời gian $t_1 = \dfrac{\Delta\theta_1}{\omega} = \dfrac{-\tfrac{\pi}{2} - (-\tfrac{2\pi}{3})}{8\pi} = \dfrac{\tfrac{\pi}{6}}{8\pi} = \dfrac{1}{48}$ s. Bước 2: Quãng đường còn lại $s_2 = s - s_1 = (2+2\sqrt{2}) - 2 = 2\sqrt{2}$ cm. Vật tiếp tục đi từ $x=0$ theo chiều dương đến vị trí $x_f = 2\sqrt{2}$ cm. Tại $x_f = 2\sqrt{2}$ cm, ta có $\cos\theta_f = \dfrac{x_f}{A} = \dfrac{2\sqrt{2}}{4} = \dfrac{\sqrt{2}}{2}$. Vì vật đang đi từ $x=0$ theo chiều dương (tức là từ pha $-\tfrac{\pi}{2}$ về $0$), nên $\theta_f = -\tfrac{\pi}{4}$. Góc quay tương ứng: từ $\theta_1 = -\tfrac{\pi}{2}$ đến $\theta_f = -\tfrac{\pi}{4}$. Thời gian $t_2 = \dfrac{\Delta\theta_2}{\omega} = \dfrac{-\tfrac{\pi}{4} - (-\tfrac{\pi}{2})}{8\pi} = \dfrac{\tfrac{\pi}{4}}{8\pi} = \dfrac{1}{32}$ s. Tổng thời gian: $t = t_1 + t_2 = \dfrac{1}{48} + \dfrac{1}{32} = \dfrac{2}{96} + \dfrac{3}{96} = \dfrac{5}{96}$ s.
}
\end{ex}
